\documentclass[conference]{IEEEtran}
\usepackage{cite}
\usepackage{amsmath,amssymb,amsfonts}
\usepackage{algorithmic}
\usepackage{graphicx}
\usepackage{textcomp}
\usepackage{xcolor}
\usepackage[hidelinks]{hyperref}

\title{A Lock Free, Non-blocking, In Line Processing Architecture for High Throughput and Regulatory Compliant Blockchain Trading Applications}

\author{Andrew Le Gear$^{1,2}$, Jim Buckley$^{1,3}$, Tawny Whatmore$^{2}$, and Ashish Sai$^{4}$ \\
$^{1}$School of Computer Science and Information Systems, University of Limerick \\
$^{2}$Horizon Globex \\
$^{3}$Lero: The Science Foundation Ireland Research Centre for Software \\
$^{4}$Maastricht University}

\begin{document}

\section*{Artifact Abstract}

\subsection*{Motivation}

Trading applications deployed on blockchain must reconcile conflicting requirements: high transaction throughput, strict regulatory compliance, and deterministic behavior. Public blockchains like Ethereum and Solana enforce per-block gas limits and impose variable transaction ordering constraints, limiting throughput severely. Teams building trading bridges typically choose between two problematic approaches: (1) sequential synchronous processing, which preserves correctness but delivers minimal throughput, or (2) fully concurrent designs that boost performance only at the expense of ordering guarantees and compliance complexity.

This work explores a middle ground: achieving high throughput while preserving the deterministic, auditable transaction ordering that regulators require—and doing so without compromising code clarity or system safety.

\subsection*{Contribution}

We explore this design space through five reference implementations, each representing a meaningful trade-off between simplicity and performance:

\begin{enumerate}
\item \textbf{Architecture 1 (Sequential Synchronous):} Baseline. Single-threaded, synchronous transaction processing. \textit{Throughput:} 0.25 tx/sec; \textit{Execution:} 3999.79 sec for 1000 transactions.
  
\item \textbf{Architecture 2 (Sequential Asynchronous):} Maintains strict transaction order via a single async queue, eliminating blocking I/O. \textit{Throughput:} 13.22 tx/sec; \textit{Execution:} 75.67 sec.

\item \textbf{Architecture 3 (Multi-Threaded Asynchronous):} Introduces multiple async consumers on a shared queue, bound to a single wallet to preserve ordering semantics. \textit{Throughput:} 18.04 tx/sec; \textit{Execution:} 55.42 sec.

\item \textbf{Architecture 4 (Unsynchronized Multi-Threaded):} Permits multiple wallets with unsynchronized producers and consumers, trading strict global ordering for parallelism. \textit{Throughput:} 24.16 tx/sec; \textit{Execution:} 41.39 sec.

\item \textbf{Architecture 5 (Symbol-Sharded, Lock-Free):} Proposed architecture. Partitions transactions by symbol (asset pair) into independent processing lanes, eliminating lock contention and preserving per-symbol ordering. Lock-free algorithm ensures non-blocking progress and regulatory compliance. \textit{Throughput:} 31.19 tx/sec; \textit{Execution:} 32.06 sec.
\end{enumerate}

Architecture 5 delivers a 124$\times$ throughput gain compared to the baseline while maintaining the deterministic, auditable ordering that regulators require per asset pair.

\subsection*{Artifact Contents}

\begin{itemize}
\item \textbf{Executable code:} C\# implementation (.NET 10) of all five architectures
\item \textbf{Docker containerization:} Multi-stage build with reproducible environment
\item \textbf{Benchmark harness:} 1000-transaction test suite with configurable parameters
\item \textbf{Performance data:} Raw results and statistical summaries
\item \textbf{Solidity contracts:} ERC-20 token and order book contract for blockchain integration
\item \textbf{Documentation:} Architecture rationale, data dictionary, and reproduction steps
\end{itemize}

\subsection*{Artifact Usage}

\subsubsection*{Requirements}
\begin{itemize}
\item Docker Desktop or .NET 10 SDK (Windows, Linux, macOS)
\item ~2 GB free disk space
\item 15 minutes per architecture run (or 90 minutes for full suite)
\end{itemize}

\subsubsection*{Quick Start}
\begin{verbatim}
cd blockchain-trading-architectures
docker-compose build
docker-compose run arch5-symbol-sharded
\end{verbatim}

Output: Transaction log, per-block statistics, and summary metrics (throughput, latency, block utilization).

\subsection*{Data Validation}

Artifact execution reproduces Table 1 from the published paper:

\begin{table}[h]
\centering
\small
\begin{tabular}{|l|r|r|r|r|}
\hline
\textbf{Architecture} & \textbf{Throughput} & \textbf{Time (s)} & \textbf{Avg Latency} & \textbf{Tx/Block} \\
\hline
Arch 1 (Sequential) & 0.25 tx/s & 3999.79 & 4.0 s & 1.0 \\
Arch 2 (Async) & 13.22 tx/s & 75.67 & 2.916 s & 52.63 \\
Arch 3 (Multi-Thread) & 18.04 tx/s & 55.42 & 2.959 s & 71.43 \\
Arch 4 (Unsync) & 24.16 tx/s & 41.39 & 27.211 s & 111.11 \\
Arch 5 (Lock-Free) & \textbf{31.19 tx/s} & \textbf{32.06} & \textbf{15.432 s} & \textbf{125.0} \\
\hline
\end{tabular}
\end{table}

Evaluators can compare Docker output against this table to confirm artifact fidelity.

\subsection*{Significance}

This artifact demonstrates practical design trade-offs that matter for real regulatory-compliant blockchain systems:
\begin{itemize}
\item Lock-free sharding preserves deterministic, auditable ordering while avoiding synchronization bottlenecks
\item Symbol-based partitioning maps directly to financial compliance requirements (per-asset audit trails)
\item Containerization ensures results are reproducible across development and evaluation settings
\end{itemize}

The work addresses a genuine gap: teams building trading applications that must bridge centralized systems to Ethereum, Solana, and Layer 2 networks face this exact problem. Our designs and implementations offer a practical roadmap.

\section*{Key References}

\begin{enumerate}
\item Le Gear, A., Buckley, J., Whatmore, T., \& Sai, A. (2026). A Lock Free, Non-blocking, In Line Processing Architecture for High Throughput and Regulatory Compliant Blockchain Trading Applications. \textit{Proceedings of the European Software Architecture Conference (ECSA 2026)}.
\item Herlihy, M., \& Wing, J. M. (1990). Linearizability: A correctness condition for concurrent objects. \textit{ACM Transactions on Programming Languages and Systems}, 12(3), 463--492. https://doi.org/10.1145/78969.78972
\item Lamport, L. (1978). Time, clocks, and the ordering of events in a distributed system. \textit{Communications of the ACM}, 21(7), 558--565.
\item Wood, G. (2014). Ethereum: A secure decentralised generalised transaction ledger. \textit{Ethereum Yellow Paper}.
\end{enumerate}

\end{document}
